\section{Systems of interest}
\label{sec:systems}

In the discussion to follow, we consider both uncontrolled and controlled systems. For our purposes, the \textbf{uncontrolled systems} category includes not only those systems types for which control is not available or required, but also those operating on predefined control sequences, such as industrial robots performing the same sets of operations over extended periods of time.  Also included are system types that can be considered uncontrolled within some time interval of interest (\textbf{decision horizon}). Systems in this category may have degradation modes that affect the system's performance within its expected useful life span. The rate of degradation is influenced by internal (e.g., chemical decomposition) and external factors (e.g., temperature of the operating environment). In addition to aforementioned industrial robots, examples of uncontrolled system types include bridges, buildings, electronic components, and certain types of rotating machinery, such as electrical power generators.

The \textbf{controlled systems} category covers all other system types, including dynamically controlled systems operating in uncertain environments, where the current state cannot be fully observed and non-determinism is present in control action outcomes. Degradation processes are influenced not only by the same kinds of internal and external factors as for the uncontrolled systems, but also by the control actions.

Most of the discussion to follow is applicable to both categories, although controlled systems would, naturally, benefit more from active DM/SHM. In describing the systems, we adopt the notation from the field of decision making under uncertainty \cite{Kaelbling1998}. 

A system \textbf{state} $s$ can be a scalar or a vector belonging to a state space $S$. A system \textbf{action} $a$ initiates state transitions, with $A$ denoting the space of all available actions ($A$ may be state-dependent). A \textbf{transition model} $T(s,a,s')$ describes the probability of transitioning to a particular state $s' \in S$ as a result of taking action $a$ from state $s$. A \textbf{reward model} $R(s,a)$ describes a positive reward obtained or a negative cost incurred as a result of taking action $a$ from state $s$. \textbf{Terminal states} form a subset $S_T \subset S$. Terminal states $S_T$ may include both failure and goal states. Transitions from a terminal state are only allowed back to itself.

If there is state uncertainty, \textbf{belief states} are used instead of regular states (also referred to as \textbf{beliefs}). A belief $b$ is a probability distribution over $S$, with $B$ denoting the space of all beliefs. \textbf{Observations} (\textit{e.g.}, sensor readings) can help with \textbf{belief estimation} and \textbf{updating}. Like a state, an observation can be a vector quantity. An \textbf{observation model} $O(s',a,o)$ describes the probability of getting an observation $o$ upon transition to state $s'$ as a result of action $a$.

The general function of decision making is to select actions. While in some systems we are only concerned with selecting a single $a_t$ at a given time $t$, decision making problems often involve selecting a sequence of actions. In a strictly deterministic system, an entire sequence of actions (a \textbf{plan}) can be selected ahead of time. In systems with action outcome uncertainty, however, a fixed plan can quickly become obsolete. Instead, a \textbf{policy} $\pi(s): S \rightarrow A$ needs to be selected that prescribes which action should be taken in any state. If a policy is selected that, in expectation, optimizes a desired metric (\textit{e.g.}, maximizes cumulative reward), it is referred to as an \textbf{optimal policy} and denoted $\pi^*$.

Throughout the paper, we use a robotic exploration rover operating on the surface of the Moon as a running example of a complex controlled system. The rover is solar-powered and stores electrical energy in a rechargeable battery.
